% =====================================================================
% Standalone note: specialisation surfaces for K3 x E and CHL BKM weights
% =====================================================================

\providecommand{\kBKM}{\kappa_{\mathrm{BKM}}}
\providecommand{\kch}{\kappa_{\mathrm{ch}}}
\providecommand{\kcat}{\kappa_{\mathrm{cat}}}
\providecommand{\kfib}{\kappa_{\mathrm{fiber}}}
\providecommand{\kHeis}{\kappa_{\mathrm{ch}}^{\mathrm{Heis}}}
\providecommand{\PhiFA}{\Phi^{\mathrm{FA}}}
\providecommand{\SpCh}{\mathrm{Sp}_{\mathrm{ch}}}
\providecommand{\Perf}{\mathrm{Perf}}
\providecommand{\Coh}{\mathrm{Coh}}
\providecommand{\Rep}{\mathrm{Rep}}
\providecommand{\Aut}{\mathrm{Aut}}
\providecommand{\Fix}{\mathrm{Fix}}
\providecommand{\NS}{\mathrm{NS}}
\providecommand{\Pic}{\mathrm{Pic}}
\providecommand{\Hdg}{\mathrm{Hdg}}
\providecommand{\CHL}{\mathrm{CHL}}
\providecommand{\Bor}{\mathrm{Bor}}
\providecommand{\Eone}{E_1}
\providecommand{\Etwo}{E_2}
\providecommand{\cT}{\mathcal T}
\providecommand{\cI}{\mathcal I}
\providecommand{\cA}{\mathcal A}
\providecommand{\cM}{\mathcal M}
\providecommand{\HH}{\mathfrak H}
\providecommand{\PP}{\mathbb P}
\providecommand{\Borch}{\mathrm{Borch}}
\providecommand{\Schur}{\mathrm{Schur}}
\providecommand{\fg}{\mathfrak g}
\providecommand{\su}{\mathfrak{su}}

\section*{Specialisation surfaces for \(K3 \times E\)}
\label{wn:sec:sigma2-moduli-k3xe}

Let \(S\) be a projective K3 surface, let \(E\) be an elliptic curve, and set
\(X=S\times E\).  A \(d=3\) specialisation datum for
\(\PhiFA_3(\Perf(X))\) consists of a complex surface
\(\Sigma_2\subset X\) together with a reference curve \(C\), here
\(C=E\), on which the residual chiral algebra lives.  The subscript in
\(\Sigma_2\) records complex dimension; it is unrelated to the
class-\(\mathcal S\) notation \(\Sigma_{g,n}\) for a genus-\(g\) curve
with \(n\) punctures.  The output is an
\(\Eone\)-chiral algebra on \(E\):
\[
 \SpCh_{\Sigma_2,E}\bigl(\PhiFA_3(\Perf(X))\bigr)\in \Eone\text{-}\mathrm{ChirAlg}(E).
\]
Any \(\Etwo\)-level braiding belongs to a named centre, double, module,
or enhancement layer such as \(\mathcal Z(\Rep^{\Eone}(-))\).  The
native \(d=3\) Stage-\(2\) output remains \(\Eone\)-chiral.

The surface classes live in \(H_4(X,\mathbb Z)\), equivalently in
\(\Hdg^{1,1}(X,\mathbb Z)\) by Poincar\'e duality.  Since \(H^1(S)=0\),
K\"unneth gives
\[
 H^2(X,\mathbb Z)=H^2(S,\mathbb Z)\oplus H^2(E,\mathbb Z),\qquad
 \NS(X)=\NS(S)\oplus \mathbb Z\,[S\times\{p\}].
\]
Thus
\[
 \mathrm{rk}\,\NS(X)=\rho(S)+1,\qquad
 1\leq \rho(S)\leq 20.
\]
The Picard-maximal algebraic surface lattice has rank \(21\).  Its
standard generators are the K3 fibre \(S\times\{p\}\) and the
surfaces \(D\times E\) with \(D\in\NS(S)\) represented by an effective
curve on \(S\).  Codimension-\(2\) curve classes lie in \(H^4(X)\);
they belong to a different lattice, outside the K3-fibre
specialisation data.

\section*{The principal K3 surface}
\label{wn:sec:principal-k3-surface}

The principal surface is
\[
 \Sigma_2^{(1)}=S\times\{p\}=p_E^{-1}(p).
\]
On the principal \(K3\times E\) recognition locus of
Theorem~\ref{thm:g-delta5-is-sp-k3}, the Stage-\(2\) algebra is
\[
 A_{S,E}:=\SpCh_{\Sigma_2^{(1)},E}
 \bigl(\PhiFA_3(\Perf(S\times E))\bigr)
 \simeq U_{\mathrm{ch}}(\mathfrak g_{\Delta_5})
\]
as an \(\Eone\)-chiral algebra on \(E\).  The automorphic test attached
to the target is
\[
 \kBKM(\Delta_5)=\mathrm{wt}(\Delta_5)
 =\frac{c_1(0)}2=\frac{10}{2}=5.
\]
The K3 elliptic-genus convention with a doubled Jacobi form belongs to
the Igusa-square normalisation.  In the theta normalisation,
\[
 [q^{1/2}r^{1/2}s^{1/2}]\Delta_5=64,\qquad
 D_5=64^{-1}\Delta_5,
\]
and the denominator algebra uses
\[
 \operatorname{den}(\mathfrak g_{\Delta_5})=D_5(2Z)
 =64^{-1}\Delta_5(2Z).
\]
The square normalisation gives
\[
 \Phi_{10}^{\theta}=(\Delta_5)^2,\qquad
 \mathrm{wt}(\Phi_{10}^{\theta})=10.
\]
Thus the primitive BKM denominator has characteristic \(5\), while its
Igusa square has weight \(10\).

\section*{Gaiotto curves, Siegel targets, and product periods}
\label{wn:sec:gaiotto-siegel-product-periods}

The class-\(\mathcal S\) parent of the Igusa corner is
\[
 \cT_{24}:=\cT[A_1,\Sigma_{0,24}],
 \qquad
 \Sigma_{0,24}=(\PP^1;p_1,\ldots,p_{24}),
\]
where all \(24\) punctures are maximal regular \(\su(2)\)-punctures.
For \(A_1\) on a punctured sphere,
\[
 n_v(\Sigma_{0,n})=3(n-3),\qquad
 n_h(\Sigma_{0,n})=4(n-2),
\]
and the Shapere--Tachikawa normalisation gives
\[
 c_{4d}=\frac{2n_v+n_h}{12},\qquad c_{2d}=-12c_{4d}.
\]
At \(n=24\),
\[
 (n_v,n_h)=(63,88),\qquad
 c_{4d}\bigl(\cT[A_1,\Sigma_{0,24}]\bigr)=\frac{107}{6},
 \qquad c_{2d}=-214.
\]

The closed genus-two class-\(\mathcal S\) curve \(\Sigma_{2,0}\) has two
trinions and three gauged \(SU(2)\) tubes, hence
\[
 (n_v,n_h)=(9,8),\qquad
 c_{4d}\bigl(\cT[A_1,\Sigma_{2,0}]\bigr)=\frac{13}{6},
 \qquad c_{2d}=-26.
\]
It has no \(24\)-puncture labels, no \(\su(2)^{24}\) flavour algebra,
and no \(M_{24}\)-permutation datum.  The Schur parent with
\((c_{4d},c_{2d})=(107/6,-214)\) is \(\cT[A_1,\Sigma_{0,24}]\).

The comparison targets have different source categories:
\[
\begin{array}{rcl}
\Sigma_2\subset S\times E
&\longmapsto&
\SpCh_{\Sigma_2,E}\bigl(\PhiFA_3(\Perf(S\times E))\bigr),\\[0.35em]
\Sigma_{0,24}
&\longmapsto&
\chi\bigl(\cT[A_1,\Sigma_{0,24}]\bigr),\\[0.35em]
Z\in\HH_2
&\longmapsto&
A_Z=\mathbb C^2/(\mathbb Z^2+Z\mathbb Z^2)\in\cA_2,\\[0.35em]
(S,E)
&\longmapsto&
[\omega_S\otimes\omega_E]\in
\PP\bigl(F^3H^3(S\times E,\mathbb C)\bigr).
\end{array}
\]
Here \(\cA_2=\mathrm{Sp}_4(\mathbb Z)\backslash\HH_2\) is the moduli
stack of principally polarised abelian surfaces, and
\[
 H^3(S\times E,\mathbb Z)=H^2(S,\mathbb Z)\otimes H^1(E,\mathbb Z)
\]
is the product variation of Hodge structure.  The Gaiotto moduli of
punctured curves, the Siegel moduli \(\cA_2\), the Humbert divisors in
\(\cA_2\), and the product period locus of \(S\times E\) meet only after
an explicit comparison datum has been supplied, such as a Kummer or
Shioda--Inose construction from a principally polarised abelian surface
to a K3 surface, or a Kuga--Satake-type period comparison with named
lattice embedding.

The Schur-to-Borcherds bridge is the typed composite
\[
 \cI_{\Schur}\bigl(\cT[A_1,\Sigma_{0,24}]\bigr)
 \xrightarrow{\ \Pi_{M_{24}}\ }
 \phi^{K3}_{0,1}(\tau,z)
 \xrightarrow{\ \Borch\ }
 \Delta_5
 \rightsquigarrow
 \fg_{\Delta_5}.
\]
The first arrow is the \(M_{24}\)-equivariant average on the
puncture-labelled Schur data; the second is the Borcherds lift of the
Gritsenko--Nikulin normalised K3 Jacobi form with \(c_1(0)=10\).

\section*{CHL correspondences and Borcherds targets}
\label{wn:sec:chl-correspondences}

For \(N\in\{1,2,3,4,6\}\), choose a symplectic automorphism
\(g_N\in\Aut(S)\) of order \(N\) and a primitive torsion translation
\(\tau_N\in E[N]\).  The diagonal action
\[
 \langle g_N\times \tau_N\rangle\curvearrowright S\times E
\]
is free because the elliptic translation has no fixed point.  The CHL
threefold is
\[
 X_N=(S\times E)/\langle g_N\times\tau_N\rangle.
\]
The cycle-level datum for \(N>1\) is therefore an equivariant
correspondence
\[
 S\times E \xleftarrow{\ q_N\ } S\times E
 \xrightarrow{\ \pi_N\ } X_N,
\]
together with an \(N\)-equivariant surface class on the quotient or on
its resolved inertia correspondence.  The fixed locus \(\Fix(g_N)\)
controls the twined Jacobi coefficient and the Du Val type of the K3
quotient; by itself it supplies fixed-point data, with complex surfaces
appearing after quotient or inertia resolution.  For nontrivial
symplectic \(g_N\), \(\Fix(g_N)\) is
zero-dimensional:
\[
 \chi(\Fix(g_N))=(8,6,4,2)\quad\text{for}\quad N=(2,3,4,6).
\]
The corresponding cyclic quotient singularity types are
\[
 8A_1,\qquad 6A_2,\qquad 4A_3+2A_1,\qquad 2A_5+2A_2+2A_1.
\]
Exceptional curves in the resolved K3 quotient and their products with
the elliptic direction provide the geometric surface pieces in the
CHL quotient geometry.

The automorphic side is already determined.  Let
\(\Delta_{k(N)}^{(N)}\) be the Gritsenko--Nikulin CHL Borcherds product
attached to the \(g_N\)-twined weak Jacobi form of weight \(0\) and
index \(1\).  The symbol \(c_N(0)\) denotes the Borcherds-input
constant in the CHL normalisation, fixed by \(c_1(0)=10\).  Borcherds'
weight formula gives
\[
 \kBKM(\Delta_{k(N)}^{(N)})=\mathrm{wt}(\Delta_{k(N)}^{(N)})
 =\frac{c_N(0)}2.
\]
The primitive CHL constants are
\[
\begin{array}{c|ccccc}
N & 1 & 2 & 3 & 4 & 6\\ \hline
c_N(0) & 10 & 8 & 6 & 4 & 2\\
\kBKM & 5 & 4 & 3 & 2 & 1
\end{array}
\]
with sources Gritsenko--Nikulin 1995 Part II, Gritsenko 1999,
Borcherds 1998 Theorem 13.3, and the CHL weight table of
Govindarajan--Krishna 2010.  These values are independent of the
compact Hodge supertrace on \(X_N\).

\begin{center}
\renewcommand{\arraystretch}{1.25}
\small
\begin{tabular}{c|l|l|c|l}
\hline
\(N\) & Geometric datum & Quotient input & \(\kBKM\) & Established input and missing construction \\
\hline
1 & \(S\times\{p\}\subset S\times E\) & identity & \(5\)
  & Theorem~\ref{thm:g-delta5-is-sp-k3} on the recognition locus \\
2 & \(\mathbb Z/2\)-CHL quotient correspondence & \(8A_1\) & \(4\)
  & Oberdieck CHL reduced-DT input; equivariant \(\SpCh\) construction required \\
3 & \(\mathbb Z/3\)-CHL quotient correspondence & \(6A_2\) & \(3\)
  & Borcherds product fixed; equivariant Hall comparison required \\
4 & \(\mathbb Z/4\)-CHL quotient correspondence & \(4A_3+2A_1\) & \(2\)
  & Borcherds product fixed; quotient-surface \(\SpCh\) datum required \\
6 & \(\mathbb Z/6\)-CHL quotient correspondence & \(2A_5+2A_2+2A_1\) & \(1\)
  & Borcherds product fixed; abelian-composition DT comparison required \\
\hline
\end{tabular}
\end{center}

Thus Corollary~\ref{cor:sibling-BKMs-from-one-phiFA} has two different
levels.  At \(N=1\) the K3-fibre specialisation is constructed under
the finite Hall--Borcherds recognition gate.  At
\(N\in\{2,3,4,6\}\) the Borcherds denominator and its weight are
theorems, while the \(\Eone\)-chiral equality
\[
 \SpCh_{\Sigma_2^{(N)},E}\bigl(\PhiFA_3(\Perf(S\times E))\bigr)
 \simeq U_{\mathrm{ch}}(\mathfrak g_{\Delta_{k(N)}^{(N)}})
\]
requires an equivariant oriented hCS-to-Hall comparison, a quotient
surface specialisation datum, PBW/no-extra-relations at the Hall
presentation, and compatibility with the Borcherds product denominator.

\section*{Gritsenko--Clery row normalisations}
\label{wn:sec:gc-normalisations}

The CHL ladder above is separate from the Gritsenko--Clery dd-modular
catalogue.  The GC classification is indexed by triples
\[
 (t,N;k)\in
 \{(1,1;5),(2,1;2),(1,2;3),(3,1;1),
   (1,3;2),(4,1;\tfrac12),(1,4;\tfrac32),(2,2;1)\}.
\]
For each row \(F_{t,N}\),
\[
 \kBKM(F_{t,N})=\mathrm{wt}(F_{t,N})=k.
\]
The doubled row constants are
\[
 2k=(10,4,6,2,4,1,3,2)
\]
in the order displayed above.  The common row with the CHL ladder is
\((t,N;k)=(1,1;5)\), where \(F_{1,1}=\Delta_5\).  Away from that row,
the group, divisor, multiplier, and input Jacobi form have their own
normalisation.  The CHL tuple \((10,8,6,4,2)\) and the GC tuple
\((10,4,6,2,4,1,3,2)\) record different moduli problems.

\section*{Humbert and Noether--Lefschetz surfaces}
\label{wn:sec:humbert-surface}

Let \(\mathcal H_D\subset\cA_2\) be the Humbert divisor of discriminant
\(D\): it parametrises principally polarised abelian surfaces whose
N\'eron--Severi lattice contains a primitive rank-two sublattice of
discriminant \(D\).  A K3 interpretation requires an additional period
bridge.  Under a Shioda--Inose/Kummer hypothesis, an abelian surface
\(A_Z\in\mathcal H_D\) gives a Kummer surface \(\mathrm{Kum}(A_Z)\) and
a K3 surface \(S\) with \(T_S\simeq T_{A_Z}(2)\).  The Humbert condition
then maps to a Noether--Lefschetz condition on \(S\): a primitive class
\(\lambda\in\NS(S)\) with the corresponding discriminant convention is
present in the K3 N\'eron--Severi lattice.

When \(\lambda\) is effective, choose a curve \(D_\lambda\subset S\).
Then
\[
 \Sigma_{2,\lambda}=D_\lambda\times E
\]
is a complex surface in \(S\times E\).  This construction supplies the
Noether--Lefschetz/Heegner surface type, separate from the CHL quotient
correspondences.

The conjectural automorphic target is the Borcherds product attached to
the Heegner divisor of \(\lambda\), as in Bruinier's singular-theta
construction.  The missing theorem is the compatibility
\[
 \SpCh_{\Sigma_{2,\lambda},E}\bigl(\PhiFA_3(\Perf(S\times E))\bigr)
 \simeq U_{\mathrm{ch}}(\mathfrak g_{\lambda})
\]
after the Heegner/Borcherds lift, the Hall presentation, and the
surface specialisation have been matched.  The required construction is
the Bruinier--Freitag singular-theta lift transported through the
\(K3\times E\) Stage-\(2\) specialisation.  The object belongs to the
Noether--Lefschetz/Heegner lane, outside the five-entry CHL ladder.

\section*{Invariant separation}
\label{wn:sec:invariant-separation}

For the compact product \(X=S\times E\),
\[
 \kcat(X)=\chi(\mathcal O_S)\chi(\mathcal O_E)=2\cdot0=0.
\]
The compact chiral Hodge supertrace also vanishes:
\[
 \kch(X)=\sum_{q=0}^3(-1)^q h^{0,q}(X)
 =1-1+1-1=0.
\]
The Heisenberg witness attached to the principal Stage-\(2\) shadow has
\[
 \kHeis(X)=3,
\]
the rank of the signature-\((2,1)\) Cartan in the
\(\Delta_5\) Borcherds comparison.  The fibre invariant is
\[
 \kfib(X)=24,
\]
the rank of the K3 Mukai lattice.  The BKM invariant at CHL level \(N\)
is
\[
 \kBKM(\Delta_{k(N)}^{(N)})=\frac{c_N(0)}2.
\]
Consequently
\[
 \{\kcat,\kch,\kHeis,\kBKM,\kfib\}(K3\times E)
 =\{0,0,3,5,24\}
\]
at the principal denominator normalisation.  The commonly quoted
\(\{0,3,5,24\}\) spectrum suppresses the compact \(\kch=0\) entry and
records the Heisenberg witness \(3\) in the four-entry shorthand.

No addition law links these columns.  At \(N=1\), any rule formed from
the compact \(\kch(X)\) and \(\chi(\mathcal O_E)\) gives \(0\), while
the Borcherds product gives \(\kBKM(\Delta_5)=5\).  At \(N=2\), the
Borcherds product gives \(c_2(0)/2=4\).  The governing identity is
\[
 \kBKM(\Delta_{k(N)}^{(N)})=c_N(0)/2.
\]

\section*{Hall and Yangian separation}
\label{wn:sec:hall-yangian-separation}

The Borcherds algebra in the principal \(K3\times E\) lane is the
Hall--Drinfeld/BKM object recognised after the compact Hall
presentation, radical quotient, and Borcherds denominator comparison.
The K3 self-mirror Yangian branch is a separate CY\(_2\) construction
from the Mukai-Heisenberg output of \(\Phi_2(D^b\Coh(K3))\).
The class-\(\mathcal S\)/VOA branch begins with
\(\cT[A_1,\Sigma_{0,24}]\) and the Schur functor; it reaches
\(\Delta_5\) through the \(M_{24}\)-averaged Jacobi form and the
Borcherds lift.  These are three branches with typed comparison maps:
Hall--Drinfeld/BKM, Yangian/Mukai, and class-\(\mathcal S\)/VOA.

The toric reference follows the same positive-half discipline:
\[
 \mathrm{CoHA}(\mathbb C^3)=Y^+(\widehat{\mathfrak{gl}}_1).
\]
The full \(\mathcal W_{1+\infty}\) algebra is recovered after the
Drinfeld double, centre, completion, or evaluation passage.  This is
the model for reading every \(d=3\) Hall comparison: the native
\(\Phi_3\) output is \(\Eone\)-chiral, and additional braided or full
vertex structure has to be named.

\section*{M5 wrapping heuristic}
\label{wn:sec:m5-wrapping}

The physical interpretation is a heuristic.  An M5-brane wraps a
four-real-dimensional specialisation surface and a circle in the
elliptic direction, with time supplying the sixth worldvolume
direction:
\[
 \mathrm{M5}\quad\text{on}\quad \Sigma_2\times S^1_E\times \mathbb R_t.
\]
For \(N=1\), \(\Sigma_2=S\times\{p\}\) gives the
\(\mathfrak g_{\Delta_5}\) BPS algebra after the Hall--Borcherds
recognition gate.  For \(N>1\), the wrapped surface lies in the CHL
quotient or its resolved inertia correspondence.  The monodromy around
\(S^1_E\) imposes the \(\mathbb Z/N\)-equivariant sector and predicts
the BPS algebra \(\mathfrak g_{\Delta_{k(N)}^{(N)}}\).  The heuristic
becomes a theorem only after the equivariant Hall presentation and
Stage-\(2\) specialisation are constructed.

\section*{Missing constructions}
\label{wn:sec:missing-constructions}

\begin{enumerate}
\item Construct the equivariant surface specialisation
\(\SpCh_{\Sigma_2^{(N)},E}\) for \(N=2,3,4,6\) from the CHL quotient
correspondence, including the resolved inertia surface classes.
\item Prove the oriented equivariant hCS-to-Hall comparison for the
same quotient geometries.
\item Match the finite Hall presentation to
\(U_{\mathrm{ch}}(\mathfrak g_{\Delta_{k(N)}^{(N)}})\) by radical
quotient, PBW, and no-extra-relations checks.
\item Transport the Bruinier--Freitag Heegner/Borcherds lift through
the Noether--Lefschetz surface \(\Sigma_{2,\lambda}=D_\lambda\times E\).
\item Keep the CHL constants, Gritsenko--Clery row constants,
Igusa-square normalisation, and the separated \(K3\times E\) invariants
in separate columns in every theorem using this comparison.
\end{enumerate}

\section*{References}
\label{wn:sec:sigma2-references}

Mukai 1988, \emph{Invent. Math.} 94, 183--221, Theorems 0.1 and 0.6,
Table 2.
Nikulin 1979, \emph{Trans. Moscow Math. Soc.} 38, Section 4.
Hashimoto 2012, \emph{Osaka J. Math.} 49, Table 5.
Gaiotto 2012, \emph{JHEP} 1208:034, arXiv:0904.2715.
Chacaltana--Distler 2010, \emph{JHEP} 1011:099, arXiv:1008.5203.
Shapere--Tachikawa 2008, \emph{JHEP} 0809:109, arXiv:0804.1957.
Beem--Lemos--Liendo--Peelaers--Rastelli--van Rees 2015,
\emph{Commun. Math. Phys.} 336, arXiv:1312.5344.
Oberdieck 2017, arXiv:1706.04948, Theorem 1.1.
Bryan--Oberdieck 2017, arXiv:1712.01593.
Bryan--Steinberg 2016, arXiv:1607.02830.
Bryan--Oberdieck--Pandharipande--Yin 2018, arXiv:1808.00846.
Gritsenko 1999, arXiv:math/9906191, Theorem 3.4.
Gritsenko--Nikulin 1995, arXiv:alg-geom/9504006, Part II.
Gritsenko--Nikulin 1998, \emph{Internat. Math. Res. Notices} 1998:15.
Gritsenko--Clery 2008, arXiv:0812.3962, Theorems 1.4 and 3.1.
Borcherds 1998, \emph{Invent. Math.} 132, Theorem 13.3.
Eguchi--Ooguri--Tachikawa 2011, \emph{Exp. Math.} 20, Table 1.
Govindarajan--Krishna 2010, \emph{JHEP} 05, 014, Table 1.
Oberdieck--Pixton 2018, \emph{Ann. of Math.} 187.
Bruinier 2002, \emph{Lecture Notes in Mathematics} 1780, Chapter 3.
Shioda--Inose 1977, \emph{Complex analysis and algebraic geometry},
119--136.
Morrison 1984, \emph{Duke Math. J.} 51, 959--976, Theorem 6.3.
Lorgat 2020, \emph{A Borcherds lift of \(\phi_{0,1}\)}, Section 2 and
Conjecture 1.
